\documentclass[ecta,nameyear,draft]{econsocart}
\usepackage{amsmath,amsfonts,amssymb,amsthm,booktabs,array,longtable}
\usepackage[T1]{fontenc}
\usepackage{mathpazo,microtype,needspace,xurl}
\usepackage{xr-hyper}
\RequirePackage[colorlinks,citecolor=blue,linkcolor=blue,urlcolor=blue]{hyperref}
\startlocaldefs
\newtheorem{theorem}{Theorem}
\newtheorem{proposition}[theorem]{Proposition}
\newtheorem{lemma}[theorem]{Lemma}
\newtheorem{corollary}[theorem]{Corollary}
\newcommand{\E}{\mathbb E}
\newcommand{\R}{\mathbb R}
\endlocaldefs
\setlength{\emergencystretch}{3em}
\makeatletter\def\copyright@text{Prepared for referee review}\makeatother

\setcounter{secnumdepth}{2}

% Stable current-document identity for clean typesetting reproduction.
\ifdefined\pdfinfoomitdate\pdfinfoomitdate=1\fi
\ifdefined\pdfsuppressptexinfo\pdfsuppressptexinfo=15\fi
\ifdefined\pdftrailerid\pdftrailerid{}\fi

\externaldocument[][nocite]{build/supp}[supp.pdf]
\externaldocument[][nocite]{build/appendix}[appendix.pdf]
\begin{document}
\begin{frontmatter}
\title{Neural Bellman Operators}
\runtitle{Neural Bellman Operators}
\begin{aug}
\author[id=au1,addressref={add1}]{\fnms{Qian}~\snm{QI}\ead[label=e92ECTA]{qiqian@pku.edu.cn}}
\address[id=add1]{Peking University}
\end{aug}
\begin{abstract}
This paper develops Neural Bellman Operators for feasible decisions under the original recursive economic objective. A centered tail-switch bound connects continuation accuracy to constrained action loss when nonlinear risk changes its maximizing density. An active-region compiler solves the original investment model's terminal Bellman action through at most four monotone cubic regions, with an outward witness retaining certification authority. Completed-incumbent safeguards guarantee actual policy improvement; recursive gain bounds and paid precision rules state the additional uniform conditions and computational costs. A Brownian BSDE result gives a separately delimited recursive-utility bridge. An executed common-information comparison retains 4,096 catalogue decisions, strong conventional representations and seven fresh-process timing repetitions. Compilation reduces observed complete work, and safeguards improve fitted policies at additional cost. Equivalent ReLU and conventional hinge geometry produces identical policies. Historical quadratic victories, adverse residual-fitting results and unresolved adoption regions remain visible. Neural-specific competitiveness and calibrated economic adoption are research claims still requiring their own evidence.
\end{abstract}
\begin{keyword}
\kwd{Dynamic programming}\kwd{Neural approximation}\kwd{Recursive risk}
\kwd{Policy qualification}\kwd{Computational economics}
\end{keyword}
\end{frontmatter}
\section{Introduction}\label{sec:intro109}
Dynamic investment under risk requires both a continuation value and a feasible decision. A small approximation error is useful only insofar as it changes the original Bellman decision at an acceptable total computational cost. Neural Bellman Operators (NBOs) address this problem by constructing continuation and policy maps, evaluating their own future decisions, and checking their losses under the original economic law. The research question is not whether a neural network can fit a value function. It is when an acquired neural continuation, compared with an equally informed conventional representation, helps deliver economically accurate dynamic decisions.

This distinction matters under nonlinear recursive risk. An approximation may improve a global residual norm while moving a nearly binding action in the wrong direction. It may also match every pointwise action derivative of the continuation outcome and yet change the risk functional's maximizing tail. Finally, a proposal that is rejected by a certificate must be repaired or replaced by a paid feasible fallback. Ignoring that last step can turn an apparent computational advantage into a false comparison between incomplete services.

We develop the original NBO construction along three linked lines. First, an original-law supporting quadratic converts continuation information, risk-witness slack and constrained stationarity into an action-loss certificate. We sharpen this connection for AVaR by localizing changes in the maximizing density to a quantile band. The resulting bound depends on centered value error, directional error and the oscillation of the action gradient in that band. An exact continuous-law counterexample shows why directional matching alone is insufficient. The finite-menu selection and information results retained below explain what can be learned or certified from the declared information; they do not make a fitted objective an economic success by definition.

Second, we construct and verify an active-region Bellman oracle for the paper's original continuous-risk investment technology. Its terminal action derivative is a strictly increasing cubic on at most four regions. This structure yields a formula-derived proposal and a state-gradient calculation, including the multiplier at a state-dependent capacity. Floating cubic inversion remains a proposal: an outward original-law witness, with the original fallback, supplies the certificate. The optimal-continuation oracle is not confused with the actual future policy whose cost is being evaluated. This is a strong semi-analytic conventional control as well as a component available to NBO; its speed is not attributed exclusively to neural representation.

Third, a completed-incumbent safeguard accepts a new action only when its verified policy-cost interval lies strictly below that of the completed incumbent. We give its exact economic interpretation, a recursive improvement theorem with lower and upper risk-distorted gain accounts, and a paid precision bound for resolving a positive improvement. A separate Brownian Markov BSDE result makes explicit the assumptions under which a neural critic and actor can be certified for a recursive-utility objective. That bridge does not identify nested AVaR with static trajectory CVaR or with an arbitrary coupled Neural Differential Utility system.

The numerical experiment executes these distinctions in the same original two-date investment economy. All fitted arms receive the same 256 paid continuation labels. Eight proposal representations include rank-six quadratic, ReLU and coordinate spline, an equivalent conventional oblique hinge, full quadratic and cubic polynomials, a tensor interpolation control, and an online original-model service. Every representation receives both the original and compiled oracle and both the plain and safeguarded complete service. We retain all 4,096 catalogue decisions and seven fresh-process timing repetitions, including refused improvements and fallback work. The catalogue is deterministic, not a new independent population sample.

The results identify real original-policy improvements and real costs. The compiler lowers observed complete-service clocks across the declared comparisons. Strict safeguarded updates reduce the fitted policies' original costs, but their additional verification is expensive relative to the already accurate online control. The same-geometry ReLU and conventional hinge give exactly identical policy and certificate records. At risk level $3/4$, their unsafeguarded catalogue mean cost is slightly below the rank-six quadratic's, but not the online control's; at $1/2$ their ordering relative to quadratic is reversed. These are numerical cost comparisons, not a neural-specific adoption claim. The historical finding that all 36 residual spans improved while aggregate rejections rose from 644 to 757 remains prominent, as do the earlier quadratic complete-work winner and the absence of certified neural adoption over it.

\paragraph{Relation to established methods.}
Risk-averse dynamic programming and policy iteration are not new in themselves \citep{ruszczynski2010risk,ruszczynski2014erratum}. Inexact-oracle optimization already separates an approximate computational step from its accuracy account \citep{DevolderGlineurNesterov2014}. Decision-focused learning and CVaR gradients likewise have established foundations \citep{ElmachtoubGrigas2022,Tamar2015}. Our mathematical contribution is the explicit connection between the original recursive objective, centered tail switching, feasible action qualification, recursive economic gain and paid resolution. Our executed contribution is a complete original-model comparison with an independently reconstructed acquisition path and a semi-analytic Bellman control. The inherited information lower bounds are stated with their exact oracle and quantifier restrictions. Neither ordinary hinge geometry nor a representation-neutral verification theorem proves neural exclusivity.

The main article keeps the original construction, current mathematical statements, full comparative evidence and supported economic interpretation together. Complete proofs appear in the Technical Supplement. Historical theorems, negative campaigns, development diagnostics and their source identities remain in the Development Appendix and source archive; preservation does not turn those separate experiments into new observations.

\section{The economic operator and its policy account}\label{sec:core49}
Consider dates $t=0,\ldots,T$ and an invariant compact state domain $X\subset\R^d$. At date $t$, the action belongs to a nonempty compact set $A_t(x)$. The transition is $F_t(x,a,Z_{t+1})$, the current cost is $c_t(x,a)$, and the terminal cost is $g(x)$. Write $B_0=1$ and $B_t=\prod_{j<t}\beta_j$, with $\beta_j>0$. An implemented policy includes its observation rule, action selection and feasible numerical realization, not just its stored parameters.

For a continuation $v$, define
\begin{equation}\label{eq:operator49}
 Q_t(v;x,a)=c_t(x,a)+\beta_t\rho_t(v(F_t(x,a,Z))),\qquad
 \mathcal T_tv(x)=\inf_{a\in A_t(x)}Q_t(v;x,a).
\end{equation}
Throughout the deterministic comparison, the evaluations $\rho_t$ are monotone and cash invariant on a common domain containing all functions and constant translates used below. They are therefore nonexpansive in the supremum norm. Continuous objectives, effective compact feasible correspondences and measurable minimizers, or measurable arbitrarily near-minimizers with their error charged, are assumed. The stipulated conditional recursion gives Bellman comparison against all admissible adapted controls. In the earlier expected-cost experiments $\rho_t$ is conditional expectation, costs are bounded and innovations have their original continuous uniform law. Section~\ref{sec:continuous95} evaluates the same investment technology under an explicitly nested AVaR preference and retains the original continuous innovation law. The direct-cost inference in Section~\ref{sec:decisions49} is specific to expectation; it is not automatically an inference theorem for nonlinear recursive utility or games.

A Neural Bellman Operator is a construction and evaluation procedure for \eqref{eq:operator49}. It forms targets using its \emph{own} fitted $f_{t+1}$, returns a neural continuation and a feasible policy, and records an economic error allowance. It does not receive the optimal continuation as a training label. Arbitrary trained networks can be evaluated by the first result below. The subsequent constructive backend specifies how a network and its actor are obtained without assuming successful nonconvex optimization.

\subsection{Action-dependent errors and complete-policy loss}
Let $H(a)$ be a true current cost and $\widehat H(a)$ its continuation approximation. If $\widehat a$ is $\delta$-optimal for $\widehat H$ on the same feasible action set, then
\begin{equation}\label{eq:center49}
 H(\widehat a)-\inf_aH(a)\le\delta+
 \operatorname{osc}_{a}\{\widehat H(a)-H(a)\}.
\end{equation}
Indeed, compare $\widehat a$ with an arbitrarily near-optimal true action and cancel the approximate objectives. The infimum limit gives the result. Additive continuation levels cancel; action-dependent discrepancies do not. Under a change in future policies, the continuation must be transported to the new economic problem before this inequality licenses reuse. The full development edition retains the original quantitative transport, recursive-domain and economic-application results. No transport bound is presumed valid after an unaccounted change of constraints or preferences.

\begin{proposition}[One-sided policy sandwich]\label{prop:sandwich49}
Suppose, uniformly in state,
\[
 f_t-\mathcal T_tf_{t+1}\le u_t,\qquad
 \mathcal T_t^\pi f_{t+1}-f_t\le v_t,
 \qquad \ell_T\le f_T-g\le u_T.
\]
Then the actual feasible policy satisfies
\begin{equation}\label{eq:sandwich49}
 0\le J_0^\pi-V_0^*\le
 \sum_{t<T}B_t(u_t+v_t)+B_T(u_T-\ell_T).
\end{equation}
The bound is invariant to date-specific constant shifts of the fitted continuations, provided the associated residual accounts are shifted consistently.
\end{proposition}
\noindent The proof is in Supplement, Section~\ref{supp:proof-core49-1}.


A sharper valid residual account for the same continuation and the same deployed actor sharpens knowledge of the policy. It does not retrain that policy or change its realized cost. Conversely, subtracting two upper bounds in \eqref{eq:sandwich49} does not identify the difference in actual costs. These distinctions govern both the historical comparisons and the new experiments.

\section{Constructing and compiling a feasible neural policy}\label{sec:construction49}
Assume the primitive objective has, for feasible state--action pairs, state modulus $A_t^Q=C_{x,t}+\beta_tF_{x,t}L_{t+1}$ and action modulus $D_t=C_{a,t}+\beta_tF_{a,t}L_{t+1}$. Let
\[
 d_H(A_t(x),A_t(z))\le K_t^A\|x-z\|_1,
 \qquad L_t=A_t^Q+D_tK_t^A,\qquad L_T=L_g.
\]
The Hausdorff condition and the primitive moduli imply that $\mathcal T_tf_{t+1}$ is $L_t$-Lipschitz whenever $f_{t+1}$ is $L_{t+1}$-Lipschitz. All moduli and numerical integration procedures are supplied by the economic primitives, not fitted from successful outcomes.

At state nodes $x_{t,i}$ with covering radius $h_t$, form effective feasible $k_t$-nets. Evaluate their own-future objectives with certified error at most $e_t$. Minimize the computed labels and retain the corresponding \emph{original} action $a_{t,i}$. A Borel repair $\Gamma_{t,i}(x)$ is feasible at $x$ and satisfies
\begin{equation}\label{eq:repair49}
 \|\Gamma_{t,i}(x)-a_{t,i}\|_1\le K_t^A\|x-x_{t,i}\|_1+\zeta_t.
\end{equation}
The allowance $\zeta_t$ concerns displacement; feasibility itself may not fail. Define
\begin{equation}\label{eq:cone49}
 f_t(x)=\min_i\{y_{t,i}+L_t\|x-x_{t,i}\|_1\},\quad
 i_t(x)=\min\arg\min_i\{y_{t,i}+L_t\|x-x_{t,i}\|_1\}.
\end{equation}
Return $\pi_t(x)=\Gamma_{t,i_t(x)}(x)$. Terminal labels approximate $g$ with error $e_T$ and use the same cone construction with modulus $L_g$.

\begin{theorem}[Own-future construction with feasible witnesses]\label{thm:construct49}
The functions in \eqref{eq:cone49} have exact affine--ReLU realizations and Lipschitz moduli $L_t$ independent of label error. The returned policy is feasible and measurable. It satisfies
\begin{align}\label{eq:enclosures49}
 -e_t\le f_t-\mathcal T_tf_{t+1}&\le D_tk_t+e_t+2L_th_t,\\
 \mathcal T_t^\pi f_{t+1}-f_t&\le e_t+D_t\zeta_t,\notag
\end{align}
and hence
\begin{equation}\label{eq:constructionloss49}
 J_0^\pi-V_0^*\le
 \sum_{t<T}B_t(D_tk_t+2e_t+2L_th_t+D_t\zeta_t)
 +B_T(2e_T+2L_gh_T).
\end{equation}
With effective covers, repairs and certified queries at the prescribed accuracy, allocating each nonnegative term a share of a positive economic tolerance gives a terminating finite construction. A fixed numerical precision or a capped finite experimental ladder need not contain that allocation.
\end{theorem}
\noindent The proof is in Supplement, Section~\ref{supp:proof-core49-2}.


The Lipschitz-extension principle is classical \citep{mcshane1934}. A native min-plus envelope and its exact ReLU circuit, using the same labels, witnesses and ties, are \emph{one method with two encodings}. They return the same function and policy. Neither Theorem~\ref{thm:construct49} nor a faster representation proves a uniquely neural capability or a training-success probability for Adam or L-BFGS. The mathematical increment is the complete feasible, own-future construction and implementation account; the empirical question is its work and policy value relative to conventional generators.


\paragraph{Primitive accuracy and resource growth.}
The complete primitive investment construction, its original-optimum accuracy theorem, and the dated horizon allocation are retained in Appendix Sections~\ref{sec:accuracy62} and~\ref{sec:complexity86}. Their state, action, innovation, and bit-growth costs remain explicit. The continuous recursive-risk investment analysis below supplies the corresponding original-law qualification used in the current comparison.


\section{Reference-centered risk kernels and complete-policy qualification}\label{sec:centeredrisk92}
A datewise worst-case error account can charge incompatible states at successive dates. Under nonlinear recursive preferences, replacing that account by an expectation under the physical transition law is not valid. We instead retain the conditional risk evaluation and center its dual representation at the actual stored NBO continuation. The resulting state-dependent certificate has an original-optimum interpretation and a finite occupation representation. It extends, rather than replaces, the signed qualification of Section~\ref{sec:risk91}.

\subsection{The centered operator}
Retain the original augmented economic state, feasible actions, nonnegative discounts, and bounded complete-policy recursions. At each $(t,x,a)$ suppose the original conditional evaluation admits the representation
\begin{equation}\label{eq:dualrisk92}
 \rho_{t,x,a}(v)=\sup_{q\in\mathcal Q_{t,x,a}}
 \left\{\int v(y)q(dy)-\psi_{t,x,a}(q)\right\}.
\end{equation}
Every $q$ is a probability measure on the original successor state space. The representation is assumed on a common bounded-function domain containing all continuations below, with the measurability and selection conditions needed for the original Bellman recursions. This is a hypothesis about the original preference primitive, not permission to replace a static risk objective by a nested one. Conditional risk duality and risk-averse dynamic programming are established principles \citep{ruszczynski2010risk,ruszczynski2014erratum}. For a nonconvex or non-cash-invariant aggregator lacking \eqref{eq:dualrisk92}, the earlier order and modulus results remain applicable.

Fix one qualified parameter $\theta$ and suppress its subscript. Its entire sequence of stored continuations is $f=(f_t)_{t=0}^T$, and its supplied complete feasible policy is $\pi$. Define, using the same $f_{t+1}$ in every term,
\begin{align}\label{eq:centered-def92}
 r_t(x,a)&=c_t(x,a)+\beta_t\rho_{t,x,a}(f_{t+1})-f_t(x),\notag\\
 k_t^f(x,a,q)&=\psi_{t,x,a}(q)+\rho_{t,x,a}(f_{t+1})
                         -\int f_{t+1}\,dq,\notag\\
 \mathcal E_{t,x,a}^f(z)&=\sup_{q\in\mathcal Q_{t,x,a}}
                   \left\{\int z\,dq-k_t^f(x,a,q)\right\}.
\end{align}
The penalty is nonnegative and has infimum zero, whether or not the supremum in \eqref{eq:dualrisk92} is attained. Direct substitution gives
\begin{equation}\label{eq:exact-center92}
 \mathcal E_{t,x,a}^f(z)
 =\rho_{t,x,a}(f_{t+1}+z)-\rho_{t,x,a}(f_{t+1}).
\end{equation}
In particular, $\mathcal E^f(0)=0$, and $\mathcal E^f$ is monotone, cash invariant, and nonexpansive in the uniform norm. Discarding $k^f$ is an upper relaxation, not this identity. Freezing one risk-maximizing kernel at $f$ is in general not even an upper relaxation.

\begin{theorem}[State-dependent original-optimum certificate]\label{thm:centeredrisk92}
Suppose verified functions satisfy $r_t^-(x,a)\le r_t(x,a)$ for every original feasible action and $r_t(x,\pi_t(x))\le r_t^{\pi,+}(x)$. Let bounded measurable $L_t,U_t$ satisfy
\begin{align}\label{eq:potential-cert92}
 L_T&\le g-f_T\le U_T,\notag\\
 L_t(x)&\le\inf_{a\in A_t(x)}
       \{r_t^-(x,a)+\beta_t\mathcal E_{t,x,a}^f(L_{t+1})\},\notag\\
 U_t(x)&\ge r_t^{\pi,+}(x)
                +\beta_t\mathcal E_{t,x,\pi_t(x)}^f(U_{t+1}).
\end{align}
Then, at every initial state,
\begin{equation}\label{eq:potential-loss92}
 f_0+L_0\le V_0^*\le V_0^\pi\le f_0+U_0,
 \qquad 0\le V_0^\pi-V_0^*\le U_0-L_0.
\end{equation}
The same statements hold at every intervening date and state. They are uniform over a qualified parameter class whenever all premises and the reported bound hold for each member with its own common continuation sequence and complete policy. The certificate permits endogenous states and actions changing at every date; it does not require a finite state space or independent deployment contexts.

With exact residuals and terminal discrepancies, equality recursions in \eqref{eq:potential-cert92} give $L_t=V_t^*-f_t$ and $U_t=V_t^\pi-f_t$. With scalar bounds from \eqref{eq:signed-input91}, their corresponding constant lower and upper sequences are feasible certificates here. Thus the exact state-dependent recursion is no looser than the signed scalar account. A strictly smaller bound can result without changing a fit, a feasible action, or an economic observation.
\end{theorem}

The last assertion concerns the best exact recursion or a verified improvement upon the constant certificates. It is not a claim that every approximately computed $L,U$ improves the scalar account. In implementation, the minimum of valid policy supersolutions and the maximum of valid optimum subsolutions retain validity, so an inferior new bound need not displace the existing one.

The lower inequality requires every original action. A discrete action search without a verified original-menu lower bound cannot establish it. For a fixed action, a feasible risk kernel supplies a lower bound on $\mathcal E^f(L)$; an upper bound requires control of the entire risk envelope. Numerical integration, feasible repairs, residual errors and penalty errors belong in these inequalities. If penalty enclosures are used, lower penalty endpoints enter an upper certificate and upper endpoints enter a lower witness. Approximate dual optimizers are accepted only through these independently checked inequalities.


\section{Continuous-law qualification in the investment model}\label{sec:continuous95}
The risk theorem becomes a numerical method only when its inequalities can be constructed from economic primitives. We now supply this construction for the original investment technology with a nested tail-risk preference. States and innovations remain continuous, every feasible action is covered, and the implemented terminal controller rather than a hypothetical exact optimizer determines policy cost. This application joins continuation construction, the original-optimum comparison and a representation comparison on one problem. It does not replace the general NBO construction or the other economic applications.

\subsection{Primitives and the object being approximated}
Let $d\geq2$ be even, $X=[0,1]^d$, and interpret subscripts cyclically. At dates $t=0,1$, the feasible investment is $0\leq a\leq A(x)=1/8+\bar x/8$, where $\bar x=d^{-1}\sum_jx_j$. Write $b_j=1/2$ for odd $j$ and $b_j=1/4$ for even $j$, and $s_j=(-1)^{j+1}$. The original transition is
\begin{equation}\label{eq:investment95}
 F_j(x,a,z)=\frac1{16}+\frac{x_j}{2}+\frac{x_{j+1}}8
       +\frac{x_j(1-x_{j+1})}{16}+b_ja+s_jz,
 \qquad z\sim\operatorname{Unif}[-1/32,1/32].
\end{equation}
Define
\begin{align}\label{eq:cost95}
 S(x)&=\frac2d\sum_j(x_j-5/8)^2+
       \frac1{4d}\sum_j(x_j-x_{j+1})^2+2(1/2-2\bar x)_+^2,\\
 g(x)&=S(x)+\frac2d\sum_j(x_j-5/8)^2,\qquad h(a)=a^2+4a^4.\nonumber
\end{align}
We use $\beta=15/16$ and, separately, $\alpha\in\{1/2,3/4\}$. Average value at risk is the upper-tail cost functional
\[
 \rho_\alpha(v)=\sup\left\{\int vq\,dP:
             0\leq q\leq(1-\alpha)^{-1},\ \int q\,dP=1\right\}.
\]
Here $P$ is the original uniform innovation law. The objective is recursively defined by $V_2=g$ and
\[
 V_t(x)=S(x)+\inf_{0\leq a\leq A(x)}
       \{h(a)+\beta\rho_\alpha(V_{t+1}(F(x,a,\cdot)))\}.
\]
It is not the static tail risk of total realized cost. A complete policy value replaces each infimum by the action prescribed by that policy, with the same recursive risk primitive.

\paragraph{Assumptions for the continuous construction.}
For this subsection, the primitives are exactly \eqref{eq:investment95}--\eqref{eq:cost95}; $\alpha$ is fixed before construction; the fitted continuation $f$ is a finite convex function on $X$; and every numerical lower or upper value is an outward enclosure. Feasible tail densities and feasible actions are retained explicitly. Terminal selection is the measurable, dyadically acquired controller described below, with its fallback included in the error bound. These concrete conditions imply compact feasible graphs, bounded costs, measurable minimizers, and the conditional-risk dual used in Theorem~\ref{thm:centeredrisk92}. They do not assert that every recursive preference or every trained network has this convex structure.

\begin{proposition}[Primitive curvature and state supports]\label{prop:curvature95}
For either risk level and every even $d$, the one-date objective
$Q_1(x,a)=S(x)+h(a)+\beta\rho_\alpha(g(F(x,a,\cdot)))$
is jointly convex on $X\times[0,1/4]$, with curvature bounded below by
\begin{equation}\label{eq:curvature95}
 \frac{4-5\beta/4}{d}\|v\|_2^2+2w^2
\end{equation}
in a direction $(v,w)$. Its value $V_1$ is convex on $X$ and satisfies
\begin{equation}\label{eq:lipschitz95}
 |V_1(x)-V_1(y)|\leq\frac{15}{d}\|x-y\|_1.
\end{equation}
Let $W$ be the objective evaluated under any fixed feasible risk density, at an anchor $(x_0,a_0)$, and let $(p_x,p_a)$ be its gradient there. For any $\lambda,\nu\geq0$, set $e=p_a+\lambda-\nu$. Then the following is a global affine minorant of $V_1$:
\begin{align}\label{eq:affine95}
 \ell(y)={}&W+\lambda(a_0-A(x_0))-\nu a_0-\frac{e^2}{4}\nonumber\\
 &+\left(p_x-\frac{\lambda}{8d}\boldsymbol{1}\right)^{\!\top}(y-x_0).
\end{align}
Neither optimal multipliers nor an exact optimizing action are required.
\end{proposition}

The curvature is obtained from the original nonlinear technology, not imposed on an interpolated value grid. The negative second-order term from $x_j(1-x_{j+1})$ is bounded using $|\partial_jg|\leq10/d$; the stage curvature dominates it. The allowance $e^2/4$ pays for imperfect stationarity. The proof in the Technical Supplement also establishes the global state modulus in \eqref{eq:lipschitz95}, including the state-dependent action constraint.

\subsection{An all-action certificate under the original innovation law}
Partition the innovation interval into $n$ equal bins $B_i$, with midpoints $z_i$ and probabilities $p_i=1/n$. An innovation bin is a set over which an inequality holds, not a replacement point mass in the economic model. Choose nonnegative $q_i\leq p_i/(1-\alpha)$ summing to one. These masses define a feasible density constant on each bin. For a fixed initial state and feasible proposal $a_0$, let $W_0$ be a lower bound on the objective under this density, obtained from globally valid affine supports of the continuation at the successor midpoints. Let $p$ enclose the derivative of the same supporting objective at $a_0$. Then
\begin{equation}\label{eq:actionlower95}
 \underline V_0(x)=\inf_{0\leq u\leq A(x)}
         \{W_0+p(u-a_0)+(u-a_0)^2\}
\end{equation}
is a lower bound on the original Bellman optimum when the supported continuation is $V_1$. The scalar quadratic is minimized over the full original interval. Interval uncertainty in $p$ is covered by taking the smaller minimum at its two endpoints, since the minimum is concave in $p$.

For a convex fitted continuation, midpoint supports give the lower risk witness and the larger endpoint value bounds the continuation throughout each bin. For $V_1$, Proposition~\ref{prop:curvature95} supplies midpoint supports and the allowance $15/(32n)$ covers each half-bin. Upper tail risk is the average of the largest $n(1-\alpha)$ upper bin bounds. This average is the exact risk of a dominating step function, hence an upper bound under the continuous law.

\begin{theorem}[Constructed original-optimum qualification]\label{thm:continuous95}
Suppose the supports and bin bounds above are outwardly verified. Suppose further that the complete terminal controller satisfies
$V_1\leq V_1^{\widehat\pi}\leq V_1+\varepsilon_1$ everywhere on $X$.
For every queried initial state, the lower value \eqref{eq:actionlower95} and the binwise upper bound $\overline J_0^{\widehat\pi}$ satisfy
\begin{equation}\label{eq:originalgap95}
 0\leq J_0^{\widehat\pi}(x)-V_0(x)
       \leq\overline J_0^{\widehat\pi}(x)-\underline V_0(x).
\end{equation}
For any stored convex continuation $f$, the same lower and upper constructions are admissible bounds for the centered-risk comparison of Theorem~\ref{thm:centeredrisk92}. Their validity does not depend on the representation of $f$, on a successful fit, or on the proposal being an exact minimizer.
\end{theorem}

To see the connection, write a verified continuation bracket as $\ell\leq V_1\leq u$ and use errors $L=\ell-f$, $U=u-f$. The centered identity gives
\[
 \rho_\alpha(f)+\mathcal E^f(L)=\rho_\alpha(\ell),\qquad
 \rho_\alpha(f)+\mathcal E^f(U)=\rho_\alpha(u).
\]
Joint evaluation avoids subtracting and then adding two separately rounded copies of $\rho_\alpha(f)$. The bound \eqref{eq:actionlower95} uses one feasible risk density to minorize the required all-action infimum; the policy upper bound controls the full risk envelope. Thus the application constructs sufficient inequalities for the original theorem. It does not claim that a non-neural continuation would receive a different validity theorem, or that computing the original-value witness is free.

The terminal controller observes each state coordinate rounded down to $40$ binary places. It accepts its floating proposal only after an outward all-action gap of at most $2^{-21}$ at that observed state; otherwise it uses rational bisection. The original-state transfer is less than $30\,2^{-40}$. Consequently $\varepsilon_1=2^{-20}$ covers the complete controller, including observation and fallback. Policy evaluation encloses this controller at every possible successor, not only at the midpoint states used by the integrator.


\section{Action-aligned acquisition and complete-policy work}\label{sec:action99}
A continuation is acquired to make a decision, not to minimize an approximation statistic in isolation. The retained investment experiment makes the distinction operational: every certificate-oriented fit reduces its validated residual span, whereas aggregate first-proposal rejections increase. We now construct an acquisition target in the units of the admission inequality. The construction uses the same original action set and risk law as the NBO policy account.

\subsection{A computable directional certificate}
Fix a state and date, and suppress their indices. Let $A\subset\mathbb R^k$ be nonempty, compact and convex, let $Q$ be the original action objective, and let $a\in A$ be the actually realized proposal. A verified supporting calculation supplies $L,U\in\mathbb R$, $s\in\mathbb R^k$ and $\mu>0$ such that
\begin{equation}\label{eq:support99}
 Q(a)\le U,\qquad
 Q(b)\ge L+s^{\mathsf T}(b-a)+\frac\mu2\|b-a\|_2^2
 \quad\text{for every }b\in A.
\end{equation}
These inequalities need not identify $s$ as a derivative of $Q$. In particular, an imperfect risk witness or continuation minorant can lower $L$ while retaining a useful supporting slope. Put $w=U-L$ and
\begin{equation}\label{eq:gamma99}
 \Gamma(a)=w-\min_{b\in A}
       \left\{s^{\mathsf T}(b-a)+\frac\mu2\|b-a\|_2^2\right\}.
\end{equation}
A certified lower bound on the displayed minimum gives a valid, possibly larger, $\Gamma$. An unverified optimizer status does not.

\begin{theorem}[Directional certificate for an NBO proposal]\label{thm:action99}
Under \eqref{eq:support99},
\[
 0\le Q(a)-\inf_{b\in A}Q(b)\le\Gamma(a).
\]
Let $\widehat s$ be a supporting slope of the surrogate used to propose $a$, let $N_A(a)=\{n:n^{\mathsf T}(b-a)\le0\ (b\in A)\}$, and suppose
$\|s-\widehat s\|_2\le d$. If a supplied $n\in N_A(a)$ satisfies
$\|\widehat s+n\|_2\le r$, then
\begin{equation}\label{eq:direction99}
 \Gamma(a)\le w+\frac{(d+r)^2}{2\mu}.
\end{equation}
For $A=[0,A_0]$, the exact minimum in \eqref{eq:gamma99} is attained at
$b=\operatorname{proj}_{[0,A_0]}(a-s/\mu)$.
\end{theorem}

The proof, including boundary actions, is in Supplement~\ref{supp:action99}. Equation~\eqref{eq:direction99} separates three trainable or computable quantities: the original-objective enclosure width $w$, the directional mismatch $d$, and the surrogate stationarity residual $r$. The residual span of the continuation is not one of these quantities. A readout can improve that span and worsen the supporting slope at a nearly binding action. Conversely, constant changes in a continuation need not affect any of them after consistent centering.

The supporting calculation is available under the original recursive-risk primitives. Suppose $F(b,\omega)=F(a,\omega)+G(\omega)(b-a)$, the stage cost is $\mu$-strongly convex, and the risk dual has an action-independent feasible witness $q$ on the original innovation space. If
\[
 V(F(b,\omega))\ge \ell(\omega)+v(\omega)^{\mathsf T}G(\omega)(b-a),
\]
then \eqref{eq:support99} holds with
\begin{equation}\label{eq:lift99}
 L=c(a)+\beta\left\{\int\ell\,dq-\psi(q)\right\},\qquad
 s=\nabla c(a)+\beta\int G^{\mathsf T}v\,dq,
\end{equation}
and any valid upper bound $U$ for the original risk objective. Integration and penalty errors must be enclosed in these quantities. The witness must remain feasible for every action being minorized; a witness feasible only at $a$ is insufficient when the risk set itself depends on the action.

In the investment economy, $G=(1/2,1/4)^{\mathsf T}$ at $d=2$, the innovation law and AVaR cap are action independent, and $h''\ge2$. Thus $\mu=2$ gives exactly the all-action quadratic of \eqref{eq:actionlower95}. The added result is the explicit link from an acquired directional support to its admission gap, rather than a new economic model or an assumption that an acquired value support is an exact derivative.


\section{An active-region compiler for the original Bellman oracle}\label{sec:compiler109}
A learned continuation must compete with the structure available in the original economy. In \eqref{eq:investment95}--\eqref{eq:cost95}, that structure includes an algebraic last-date Bellman oracle. Giving this oracle only to a neural method would confound representation with information. We derive it explicitly and give its implementation to every comparator.

Put $R=1/32$, $t=1-\alpha$, $r_j=F_j(x,0,0)$, and
\begin{align*}
 H_x&=\tfrac12-2\bar r,& k_x&=\frac{10}{d}\sum_js_jr_j,\\
 C_x&=\frac8d\sum_jb_j(r_j-5/8)
 +\frac1{2d}\sum_j(b_j-b_{j+1})(r_j-r_{j+1}).
\end{align*}
The following formula uses the original continuous law, not a fitted transition or a finite-shock replacement.
\begin{theorem}[Exact tail algebra and constrained Bellman regions]\label{thm:compiler109}
For $z\sim\operatorname{Unif}[-R,R]$ and $0<\alpha<1$, define
\begin{equation}\label{eq:tailpoly109}
 \mathcal R_\alpha(k)=\operatorname{AVaR}_\alpha(5z^2+kz)
 =\begin{cases}
 \dfrac{5R^2(1+\alpha+\alpha^2)}3+\dfrac{\alpha k^2}{20t},& |k|\le10Rt,\\[2mm]
 \dfrac{5R^2(4\alpha^2-2\alpha+1)}3+\alpha R|k|,& |k|\ge10Rt.
 \end{cases}
\end{equation}
The two expressions and their first derivatives agree at their junctions. For every even $d$ and every original state $x$, the last-date action derivative is
\begin{equation}\label{eq:cubic109}
\begin{split}
 \partial_a Q_1(x,a)={}&16a^3+2a+\beta\{C_x+41a/32-3(H_x-3a/4)_+\}\\
 &+\frac{5\beta}{4}\operatorname{clip}\!\left(\frac{\alpha(k_x+5a/4)}{10t},-\alpha R,\alpha R\right).
\end{split}
\end{equation}
It is continuous and strictly increasing, with slope at least two wherever its second derivative is defined. The interval $[0,A(x)]$ has at most four regions separated by
$H_x/(3/4)$ and the two solutions of $k_x+5a/4=\pm10Rt$.
On each region \eqref{eq:cubic109} is $16a^3+pa+q$, with $p\ge2$.
There is a unique constrained minimizer: zero or capacity when the corresponding endpoint derivative has the appropriate sign, and otherwise the unique interior root. With $u=\sqrt{p/48}$, that root on its region is
\[
 a=-2u\sinh\!\left\{\tfrac13\operatorname{arsinh}(q/(32u^3))\right\}.
\]
Coefficient preparation costs $O(d)$ arithmetic; region identification and the number of scalar inversions are bounded independently of $d$.
\end{theorem}

The complexity statement counts real arithmetic and scalar inversions, not unit-cost exact transcendental evaluation. The implementation uses floating inversion only to propose a dyadic action. It then applies the unchanged outward original-law supporting certificate, including the state-capacity multiplier and its stationarity allowance. A failed compiled certificate pays for its failed attempt and the original terminal oracle. Consequently compiler error can change work but cannot acquire certificate authority. The full terminal policy used in nested policy-cost evaluation is not replaced by the compiler. In particular, more accurate optimal-continuation information and a different deployed future policy are distinct interventions.

\paragraph{Gradients and the computational graph.}
The last-date value is continuously differentiable on the relative interior of the state cube, with the one-sided boundary interpretation, and
\begin{equation}\label{eq:envelope109}
 \nabla V_1(x)=\nabla_xQ_1(x,a^*(x))-\lambda(x)\nabla A(x),\qquad
 \lambda(x)=\boldsymbol1_{\{a^*=A(x)\}}[-\partial_aQ_1(x,a^*)]_+.
\end{equation}
Within a smooth interior action region,
$\nabla a^*=-\partial_{ax}Q_1/\partial_{aa}Q_1$; within a strictly active endpoint region it is respectively zero or $\nabla A$. No global differentiability of the actor across region changes is asserted. In particular, differentiating the rounded actor or treating a feasible witness slope as this exact gradient is incorrect. Supplement Section~\ref{supp:compiler109} proves the formulas, uniqueness, boundary rules and certificate invariance. The original general own-future construction remains Theorem~\ref{thm:construct49}; the new experiment isolates its acquired last-continuation slice, not a newly trained long-horizon network.

\section{Tail-aware action accuracy}\label{sec:tail109}
A critic can match action derivatives pointwise in the innovation and still propose the wrong risk-averse action. The missing quantity is the change in the adverse tail. This section refines the directional account in Theorem~\ref{thm:action99}; it does not replace an original-law witness by a training statistic.

Fix a state and a realized action $a$. On the same original probability space let $Z_a$ be the true continuation outcome and $H_a$ an acquired surrogate. In a neighborhood of $a$, suppose both are almost surely differentiable with integrably dominated derivatives. At $a$ assume their $\alpha$-quantiles are unique, carry no mass, and cut off upper tails of probability $t=1-\alpha$, where $0<\alpha<1$. These are local assumptions, not automatic properties of every fitted ReLU. Choose a constant $c$ and certified numbers $\epsilon,\delta\ge0$ such that
\[
 \|Z_a-H_a-c\|_\infty\le\epsilon,\qquad
 \|\nabla_aZ_a-\nabla_aH_a\|_\infty\le\delta.
\]
The constant is an anchor-level centering; no derivative of a data-dependent centering rule is taken. Write $\eta$ for the true quantile, $B_\epsilon=\{|Z_a-\eta|\le2\epsilon\}$, and, for an arbitrary vector $v_0$,
\[
 G_\epsilon=\operatorname*{ess\,sup}_{B_\epsilon}\|\nabla_aH_a-v_0\|_2,
 \qquad m_\epsilon=\min\{P(B_\epsilon),2\min(\alpha,t)\}.
\]
Set $G_\epsilon=0$ when $P(B_\epsilon)=0$.
\begin{theorem}[Centered tail-switch bound]\label{thm:tail109}
Under these hypotheses, whenever $G_\epsilon<\infty$,
\begin{equation}\label{eq:tailgradient109}
 \left\|\nabla_a\operatorname{AVaR}_\alpha(Z_a)
       -\nabla_a\operatorname{AVaR}_\alpha(H_a)\right\|_2
 \le \delta+G_\epsilon m_\epsilon/t.
\end{equation}
If $P(|Z_a-\eta|\le r)\le Cr^\kappa$ for the relevant $r$, the tail term is at most $G_\epsilon C(2\epsilon)^\kappa/t$. If additionally
$\|\nabla_aH_a-v_0\|_2\le L|Z_a-\eta|$ on $B_\epsilon$, it is at most
$LC(2\epsilon)^{\kappa+1}/t$.

Suppose the original current objective $Q$ is differentiable and $\mu$-strongly convex on a compact convex action set $A$, and its stage cost is unchanged in the surrogate $\widehat Q$. If $n\in N_A(a)$ and $\|\nabla\widehat Q(a)+n\|_2\le r_a$, then
\begin{equation}\label{eq:tailaction109}
 0\le Q(a)-\min_AQ
 \le \frac{\{r_a+\beta(\delta+G_\epsilon m_\epsilon/t)\}^2}{2\mu}.
\end{equation}
\end{theorem}

The proof localizes changed tail membership to a quantile band and uses the zero total mass of the difference between two maximizing densities. Thus a constant component of the surrogate action gradient cancels: tail \emph{gradient oscillation}, rather than its overall magnitude, controls the switching cost. Under a positive margin exponent, joint control of the centered value error and directional error gives a rate for action loss. A finite validation sample does not itself verify $C$, $\kappa$ or the essential suprema. At atoms or nonunique tails the stated differentiability theorem is unavailable; the feasible-density minorants of Theorem~\ref{thm:action99} remain the appropriate certification device. When $Z_a$ uses an optimal continuation, \eqref{eq:tailaction109} bounds the current-action gap; an imperfect implemented future policy requires its separate recursive allowance before interpreting the bound as complete-policy regret. Numerical witness slack is paid there through $U-L$, not silently dropped from \eqref{eq:tailaction109}.

\paragraph{An exact diagnostic.}
Let $z$ be uniform on $[-1,1]$, $\alpha=1/2$, $a\in[-1,1]$, and stage cost $a^2$. Take $Z_a=(a+1)z$ and $H_a=(a-1)z$. Their action derivatives are identically $z$, yet the original and surrogate optimal actions are $-1/4$ and $1/4$. Using the surrogate action incurs original loss $1/4$. This continuous-law diagnostic, independently integrated in rational arithmetic, rules out any general action-loss bound that vanishes with directional error alone. It is not the investment economy, a new population sample, or an empirical neural advantage. Its conceptual provenance in the earlier tail-aware research protocols is retained.

\paragraph{Acquisition target and novelty.}
For a family satisfying verified margin assumptions, a useful acquisition target is the bound in \eqref{eq:tailaction109}, with envelope slack and actual qualification work added. The earlier prospective finite-menu selection theorem then applies to complete services with that verified certificate. Our executed comparison below does not retroactively become a training experiment for this new loss: its frozen acquisition remains least squares. CVaR gradient estimation and decision-focused learning already have established formulations \citep{Tamar2015,LuoDelage2026,ElmachtoubGrigas2022}. The present result is a centered, tail-localized perturbation bound tied to the same constrained Bellman action and a valid complete-policy account. The static trajectory-CVaR objective of \citet{LuoDelage2026} is not the nested AVaR preference used here; its empirical gains cannot be transferred between those objectives.

\section{Safeguarded policy improvement and resolution}\label{sec:improvement109}
Qualification and improvement answer different questions. A completed policy may satisfy a loss tolerance without minimizing cost, and a rejected early proposal is not the completed conventional baseline. We retain a qualified incumbent before attempting any improvement.

\begin{proposition}[Completed-incumbent improvement]\label{prop:safe109}
Fix the original state, law, feasible action set and complete future policy. Suppose two feasible current actions have verified complete-policy cost intervals $[l,u]$ and $[l',u']$. Replacing the incumbent only when $u'<l$ guarantees a strict cost decrease of at least $l-u'>0$. Repeating this rule preserves weak improvement relative to the initial \emph{completed} incumbent. Requiring every replacement also to pass the original-optimum and policy-width targets preserves those targets. A refused trial leaves the incumbent unchanged but still incurs its full work.
\end{proposition}

This rule is an interval implementation of a descent safeguard, not a new claim that descent is uniquely neural. In the executed service, three projected support steps are attempted, with fractions $1,1/2,1/4$, after the original fallback and refinement have completed. An original-optimum lower bound can be reused across actions at the same state. A policy-cost interval cannot. The projection generates candidates; only the outward gate grants a policy change. Both learned and online services receive exactly this option.

\subsection{When local improvement becomes recursive improvement}
A finite catalogue does not establish a statewise policy theorem. To state the required condition, consider bounded measurable functions, bounded costs, $0<\beta<1$, and original fixed-policy operators
\[
 T_\pi v(x)=c(x,\pi(x))+\beta\rho_{x,\pi(x)}(v(F(x,\pi(x),Z))).
\]
Assume the operators preserve bounded measurability, the one-step risks are monotone and cash invariant, and the Bellman infimum admits measurable minimizers or effective near-minimizers with their errors charged. Assume also a common dual representation, at each fixed $(x,a)$,
$\rho_{x,a}(v)=\sup_{q\in\mathcal Q(x,a)}\{\int v\,dq-\psi_{x,a}(q)\}$,
where each $q$ is a probability law. The dual set may depend on the action, but the two continuations compared at that action use the same set and penalty.

\begin{theorem}[Recursive improvement with an economic gain account]\label{thm:improve109}
The original Bellman operator $T=\inf_\pi T_\pi$ and every $T_\pi$ are $\beta$-contractions. They have unique bounded fixed points $V^*$ and $J^\pi$. Suppose the \emph{uniform} improvement residual satisfies
\[
 0\le d(x)\le J^\pi(x)-T_{\pi'}J^\pi(x)\le e(x)
\]
for bounded nonnegative $d,e$. Let $P^-v=\inf_{q\in\mathcal Q(x,\pi'(x))}\int v(F)\,dq$ and let $P^+$ replace the infimum by a supremum; assume these maps preserve bounded measurability. Their unique discounted fixed points
\[
 D^-=d+\beta P^-D^-,\qquad D^+=e+\beta P^+D^+
\]
satisfy
\begin{equation}\label{eq:gain109}
 0\le D^-\le J^\pi-J^{\pi'}\le D^+.
\end{equation}
These gain accounts use risk-distorted future laws, not physical state frequencies.

If a sequence of such uniformly safe updates additionally satisfies
$T_{\pi_{k+1}}J^{\pi_k}\le TJ^{\pi_k}+\varepsilon_k$, then
\begin{equation}\label{eq:policyiteration109}
 \|J^{\pi_{k+1}}-V^*\|_\infty
 \le\beta\|J^{\pi_k}-V^*\|_\infty+\varepsilon_k.
\end{equation}
In particular $\varepsilon_k\to0$ implies convergence to the original optimum. A verified bound $0\le J^\pi-TJ^\pi\le(1-\beta)\varepsilon$ is a valid stopping certificate for $\|J^\pi-V^*\|_\infty\le\varepsilon$.
\end{theorem}

Risk-sensitive policy iteration is established theory \citep{ruszczynski2010risk,ruszczynski2014erratum}. Here \eqref{eq:gain109} states which original-objective residual an outward improvement gate must establish and how its economic gains propagate. The finite executed root catalogue verifies Proposition~\ref{prop:safe109}, not the uniform premises of Theorem~\ref{thm:improve109}. Neither three trial steps nor successful critic fitting establishes global policy-iteration convergence.

\subsection{Paid precision and a stopping rule}
A strict descent gate can be undecided because two policy intervals overlap. Suppose a fixed incumbent and candidate have actual improvement $\Delta=J-J'$, and refinement at resolution $n$ gives intervals containing their respective costs, with widths at most $A/n^2$ and $A'/n^2$. Then
\begin{equation}\label{eq:precisiongate109}
 l_n-u'_n\ge\Delta-(A+A')/n^2.
\end{equation}
For $\Delta>0$, doubling $n$ until $n>\sqrt{(A+A')/\Delta}$ certifies improvement. With $n_0=1$ and cost at most $cn$ per refinement round, the total paid work is less than
$4c\max\{1,\sqrt{(A+A')/\Delta}\}$, apart from fixed setup. If the proposed action has loss and width targets too, those must be verified separately. If $\Delta=0$, no finite strict-descent certificate is promised. Moreover, from interval information alone, overlap is an information obstruction: an algorithm cannot certify strict descent for every pair of costs consistent with two overlapping intervals. Thus the claimed resolution scale belongs to this interval interface, not a new universal lower bound for all economic oracles.

For an operational stopping decision, let $b\le V^*(x)$ be valid and let $u$ bound the incumbent cost. Its remaining possible cost reduction is at most $u-b$. If a known nonnegative compute price converts the unavoidable expense of a contemplated improvement operation to at least $u-b$ in the same units, that operation cannot provide strictly positive net gain at this single decision. For a reused policy, multiply the qualified gain by its declared workload before comparing with acquisition expense. This dominance rule is implementable once costs and prices are supplied; it is neither an empirical adoption claim nor a globally optimal search policy over unknown future discoveries. Supplement Section~\ref{supp:improve109} gives complete proofs.

\paragraph{Recursive utility, information and the computational graph.}
Supplement Section~\ref{supp:bridge109} proves a separate Brownian Markov BSDE verification result, including well-posedness, terminal or boundary errors and the critic graph $z=\sigma^{\mathsf T}\nabla u$. Its assumptions delimit a rigorous subcase of a connection to Neural Differential Utility; it is not an equivalence between different recursive objectives. The following information results retain their own oracle restrictions and do not assume free formula or derivative access.

\section{Decision-scale numerical and statistical precision}\label{sec:precision99}
The continuous qualification uses the original terminal-controller allowance and all-action lower witness, not just a small surrogate residual. Convex midpoint and endpoint risk rules supply a width of at most $KR/\{(1-\alpha)n^2\}$ with $2n+1$ continuation queries. Theorem~\ref{thm:nonuniform99} below extends this account beyond the uniform law. The full uniform-rule statement, numerical endpoint charges and tensor extension remain in Section~\ref{sec:precision96} of the Development Appendix.

For the frozen population comparison, every complete service is qualified to the original-optimum allowance $\tau=1/256$. The design uses 4,096 independent initial states at each risk level and twelve simultaneous contrasts, with numerical half-width and sampling allowance totaling $3/8192$. Theorem~\ref{thm:power96} states the exact simultaneous event and its independence requirements. Those population intervals belong to the frozen actors, not to subsequently refitted candidates. The acquisition comparison instead retains the deterministic contrast support $[-\tau,\tau]$.

For a workload of $B$ decisions, the purchase account combines the per-decision policy contrast, $B$ times that contrast, full-workload computation and one-time acquisition. Theorems~\ref{thm:purchase97} and Proposition~\ref{prop:purchase96} retain the complete interval and strict-boundary statements. Historical timing ranges are descriptive inputs to a conditional price exercise. The exact decision mechanism in Section~\ref{sec:decision99} will provide an identified model-specific policy contrast; it will not turn these historical clock ranges into confidence intervals or external prices into estimated welfare weights.

\section{Information resolution under nonuniform laws and randomization}\label{sec:information99}
The information required to qualify a continuation is distinct from the information supplied when its formula is already known. This section studies point-value access to a convex, $K$-Lipschitz continuation on $[-R,R]$, with $K,R>0$. The continuation formula, derivatives and exact Bellman labels are not supplied for free. The original uniform-law result, including its complete proof and constants, is retained as Theorem~\ref{thm:information97} in the Development Appendix. We extend the law and distinguish two nonequivalent meanings of randomized certification.

\subsection{A known nonuniform innovation law}
Let $P$ have known density $p$ satisfying $0<m\le p(z)\le M_p<\infty$ on $[-R,R]$. Fix $0\le\alpha<1$, put $t=1-\alpha$, let $q$ be the $\alpha$-quantile of $P$ (with $q=-R$ at $\alpha=0$), and set $L=R-q$. Let $\mathcal C_{K,R}$ denote the convex $K$-Lipschitz functions on this interval.

\begin{theorem}[Nonuniform-law information and enclosure]\label{thm:nonuniform99}
Every deterministic adaptive procedure using at most $N$ point values and reporting an interval valid for all consistent $f\in\mathcal C_{K,R}$ has an input with reported width at least
\begin{equation}\label{eq:lower99}
 \frac{mKL^2}{12t(N+1)^2}
 \ \ge\ \frac{mKt}{12M_p^2(N+1)^2}.
\end{equation}
For an upper bound, partition $[-R,R]$ into $n$ equal cells. Suppose the original cell masses and conditional cell means are available exactly. The mass-weighted conditional-mean rule and the mass-weighted endpoint secant rule bracket $\operatorname{AVaR}_\alpha^P(f)$ and have width at most
\begin{equation}\label{eq:upper99}
 \frac{2M_pKR^2}{t n^2}.
\end{equation}
They use at most $2n+1$ distinct continuation values. Thus the deterministic error-versus-information rate is $N^{-2}$ for each fixed admissible law and fixed $\alpha$. Certified, rather than exact, law moments require their additional numerical allowances.
\end{theorem}

The lower bound uses indistinguishable convex tails under the original density, not a probability transformation that might destroy convexity. The upper bound applies conditional Jensen and secant inequalities to $(f-\eta)_+$ before minimizing over the AVaR threshold $\eta$. The moments are law information and their acquisition cost is explicit. With $m=M_p=1/(2R)$ and $L=2Rt$, both constants reduce to their uniform-law counterparts. Their opposite behavior as $\alpha$ tends to one remains; no tail-uniform constant is identified.

\subsection{Sound random intervals and randomized point estimates}
A randomized algorithm can mean a zero-failure enclosure or a statistical approximation. We call an interval procedure \emph{pathwise sound} if, outside one null set of random tapes independent of the input, it contains the risk of every admissible input consistent with the tape's transcript. This is stronger than requiring coverage $1-\delta$ separately for each input.

\begin{proposition}[Randomized information with explicit quantifiers]\label{prop:random99}
The first lower bound in \eqref{eq:lower99} holds on one fixed quadratic-tail input for every sound random tape of a pathwise-sound interval procedure. It therefore also bounds expected width on that input.

For ordinary randomized point-valued procedures under the uniform law, let $A(f)$ use at most $N\ge1$ point values on every realization. Then
\begin{equation}\label{eq:randomlower99}
 \sup_{f\in\mathcal C_{K,R}}
 \left\{\mathbb E\left|A(f)-\operatorname{AVaR}_\alpha^P(f)\right|^2\right\}^{1/2}
 \ge\frac{KR(1-\alpha)}{960\,N^{5/2}}.
\end{equation}
The expectation is over the algorithm's randomization. This is a lower bound, not a claim that $N^{-5/2}$ is the optimal randomized rate on the whole class.
\end{proposition}

For the first assertion, the adverse quadratic is fixed before the random tape; its indistinguishable chord continuation can depend on the realized transcript because soundness is simultaneous over inputs. That reasoning cannot be applied to a merely high-probability interval. The second assertion instead uses a fixed finite distribution of convex perturbations and the conditional variance of unrevealed signs. Its proof supplies all constants and treats adaptive queries. Classical information-based complexity already uses indistinguishable inputs and randomized bump distributions \citep{Novak2015Integration97}; the result here specializes those tools to the convex monotone AVaR tail and states which certification quantifier survives randomization.

\subsection{Dated resolution and the Bellman resource account}
Suppose, on a fixed acquired state--action graph, date $j$ contributes a verified nonnegative allowance $a_j/n_j^2$ to the original-policy tolerance and incurs $c_j n_j$ continuation-query work, with $a_j,c_j>0$. The coefficients include discounting, the appropriate tail-risk factor and the number of original objectives to be evaluated. The continuous allocation problem has value
\begin{equation}\label{eq:allocation99}
 \min_{n_j>0:\,\sum_j a_j/n_j^2\le\varepsilon}
       \sum_j c_jn_j
 =\frac{\left(\sum_j a_j^{1/3}c_j^{2/3}\right)^{3/2}}{\sqrt\varepsilon}.
\end{equation}
Rounding each minimizing $n_j$ upward preserves the tolerance and adds less than $\sum_jc_j$ query work. This is a constructive composition of certified error accounts, not an end-to-end Bellman lower bound. Acquiring the state--action graph, the neural continuation, law moments, exact arithmetic and feasible owners remains additional work. With multiple innovations a direct tensor implementation still has product point counts. The original horizon and vector-control results in Proposition~\ref{prop:horizon95} and the Development Appendix are retained, including their adverse executed scaling evidence.

\section{A decision-theoretic value of acquired oblique geometry}\label{sec:decision99}
An approximation separation matters economically only when it changes an implemented decision. We give such a realization for the original oblique continuation, without altering the retained investment comparison. This is an exactly solved economic mechanism, not an empirical calibration or a timing experiment.

A decision maker observes $x=(x_1,x_2)\in[0,1]^2$ and a posted price $p\in[0,3]$. Purchasing one durable claim yields next-date benefit
\begin{equation}\label{eq:benefit99}
 v(x)=1+(x_1+x_2-1)_+.
\end{equation}
The feasible action is $b\in\{0,1\}$, and total cost is $b\{p-v(x)\}$. This is a two-date Bellman problem with a deterministic terminal payoff and unit discount. A stored continuation $h(x)\in[0,3]$ implements $b_h(x,p)=\boldsymbol1\{p\le h(x)\}$. The original optimum buys at $p\le v(x)$. A price tie has zero loss. Let states be uniform on the square and prices uniform on $[0,3]$, independently.

\begin{theorem}[Exact policy-loss separation]\label{thm:decision99}
For every measurable $h:[0,1]^2\to[0,3]$, the conditional and unconditional policy regrets are
\begin{equation}\label{eq:decisionloss99}
 \mathbb E_p[J^{b_h}(x,p)-J^*(x,p)]
       =\frac{(h(x)-v(x))^2}{6},\qquad
 \mathbb E[J^{b_h}-J^*]=\frac{\|h-v\|_{L^2}^2}{6}.
\end{equation}
Among arbitrary square-integrable additive continuations and among polynomials of total degree at most two, respectively, the optimal expected regrets are
\begin{equation}\label{eq:separation99}
 \frac1{540}\quad\text{and}\quad\frac1{5400}.
\end{equation}
They are attained by the admissible functions
\begin{align}
 h_A(x)&=\frac56+\frac{x_1^2+x_2^2}{2},\label{eq:additive99}\\
 h_Q(x)&=\frac{17}{15}-\frac35(x_1+x_2)
              +\frac12(x_1^2+x_2^2)+\frac65x_1x_2.\label{eq:quadratic99}
\end{align}
A single oblique ReLU with the affine constant in \eqref{eq:benefit99} attains zero regret. A conventional implementation of the identical oblique hinge also attains zero regret and the same decisions.
\end{theorem}

The shift by one is part of the economic construction. It keeps the projection optima inside the posted-price support, so clipping does not secretly enlarge a comparator class. The proof computes the additive conditional-expectation projection and the six-dimensional polynomial Gram system exactly. It establishes a policy comparison against the best member of each specified class, rather than comparing arbitrarily chosen fits. The scalar price integration is the step that turns continuation error into action loss; no such identity is presumed in the investment economy.

The sup-norm comparison also has a decision interpretation. Under the same conditional price law, the smallest worst-state expected regrets for the two classes are $1/96$ and $1/1536$, respectively. These follow from the sharp errors $1/4$ and $1/16$ in Proposition~\ref{prop:ridge97}, with its explicit feasible attaining functions shifted by one. The population optima in \eqref{eq:additive99}--\eqref{eq:quadratic99} are different from the minimax optima. An acquisition objective must specify which context criterion it is intended to serve.

\subsection{Acquisition, conventional geometry and purchase}
The mechanism identifies a value of the oblique continuation geometry, not a value of relabeling its formula as a network. A piecewise-affine conventional representation on the two triangles separated by $x_1+x_2=1$ is also exact. Convex regression with that hinge already in its dictionary contains the same exact continuation. Those controls are included in the exact comparison. Discovering the separating hyperplane from unknown data, fitting a larger dictionary or constructing an adaptive triangulation are different acquisition problems; their costs cannot be inferred from this formula-access calculation.

Suppose the exact oblique service has additional fixed acquisition cost $K\ge0$, incremental per-decision computation cost $c$, and fee $F\ge0$, all expressed in the units of this model's cost. Against the best quadratic continuation, the exact gain from $m$ independent decisions is
\begin{equation}\label{eq:exactpurchase99}
 G_Q(m)=m(1/5400-c)-K-F.
\end{equation}
When $c<1/5400$, the first strictly beneficial workload is
$\lfloor(K+F)/(1/5400-c)\rfloor+1$. This is a nonzero, exactly identified policy benefit with an explicit acquisition threshold. It is not an estimate of the price of computation in the retained investment application. Equal hinge and neural implementations share this gain when their complete costs agree.

\begin{table}[htbp]
\caption{Exact original-law policy loss in the oblique-continuation mechanism}\label{tab:decision99}
\centering\small
\begin{tabular}{@{}p{.52\linewidth}p{.17\linewidth}p{.22\linewidth}@{}}
\toprule
Continuation class or supplied geometry & Optimal mean regret & Information boundary\\\midrule
Arbitrary square-integrable additive functions & $1/540$ & Population projection\\
Total-degree-two polynomials & $1/5400$ & Six exact moments\\
One oblique ReLU and affine constant & $0$ & Supplied hinge\\
Identical conventional oblique hinge & $0$ & Same supplied hinge\\
Piecewise affine on the two separating triangles & $0$ & Same separating line\\
Convex dictionary containing the oblique hinge & $0$ & Contains exact function\\\bottomrule
\end{tabular}
\par\smallskip\parbox{.96\linewidth}{\footnotesize Notes: States are uniform on $[0,1]^2$ and independent prices uniform on $[0,3]$. Loss is evaluated against the original purchase optimum. The last four rows are equivalent exact geometric realizations, not four independent fitted experiments. The quadratic row uses six right-side moments and its known Gram matrix; moment acquisition is not free when the continuation formula is unknown. No timing, empirical calibration or dimension-wide neural dominance is inferred.}
\end{table}

The table is generated from two separately implemented exact-arithmetic calculations. The production path assembles monomial moments and solves normal equations by rational elimination. The independent path integrates polynomials after a triangular change of variables and uses an orthogonal shifted-Legendre basis. It imports neither the production integrator nor its solver. Both reproduce the fitted function, the original-law loss and the range checks. This independently checks acquisition for the stated finite-dimensional mechanism; it does not independently refit the R97 nonlinear dictionary search or complete the separately frozen R98 campaign.



\section{Original-economy comparison: policies, complete work and limits}\label{sec:study109}
\subsection{Design, information and execution identity}
The experiment reproduces the executable source frozen at commit \texttt{29180a4} under its original R108 protocol, with a distinct R109 execution identity. It does not complete the separately frozen R98 or R100--R107 campaigns. Economic primitives, two dates, two state coordinates, scalar investment capacity, discount $15/16$, continuous uniform innovation and nested AVaR levels $1/2$ and $3/4$ are unchanged. All completed policies meet an original-optimum loss bound of $1/256$ and a policy-cost interval-width bound of $1/8192$. The actual terminal controller and its $2^{-20}$ allowance are unchanged.

Training uses every midpoint of a $16\times16$ grid, with paid original-model terminal-value enclosures as labels. Evaluation uses all 64 midpoints of a disjoint $8\times8$ grid. The grids were fixed before the frozen campaign, but overlap catalogues proposed in earlier research; they are not an unseen statistical sample. All means below refer to this complete equal-weight finite catalogue. The intervals are outward numerical bounds, not confidence intervals for a continuous initial-state population. R96 or R97 statistical intervals are never assigned to these new fits.

The rank-six quadratic, ReLU and coordinate-spline arms use the same 32-feature dictionary sizes, three greedy nonlinear selections, nonnegative nonlinear coefficients, specified numerical tie rule and 30-bit readout rounding. A conventional hinge arm evaluates exactly the acquired ReLU geometry with different code. Full degree-two and degree-three polynomial fits and bilinear grid interpolation broaden the conventional resource frontier; their unequal storage is recorded rather than mislabeled memory matching. All arms use identical labels and the same original-law qualification, repair and fallback rules. Full polynomial and grid proposal optimization uses the frozen 33-point scan and adjacent bounded scalar search, not a free exact policy.

Acquisition is independently reconstructed using coordinate-wise feature assembly, nonlinear active-set enumeration and pivoted QR, compared with the production constrained solver and SVD. Every nonexcluded feature-trial objective is compared, and all six rank-six feature choices and rounded coefficient vectors agree exactly in this execution. Numerical libraries are partly shared; this is algorithmic replication, not exact optimization certification or independent human authorship. Every serialized model is reloaded and its training predictions checked. ReLU and hinge share and pay for the same mathematical acquisition and are not counted as two independent learning successes.

\subsection{Complete services and economic results}
Each arm is evaluated with the legacy oracle, the compiled oracle, and a safeguarded version of each. The starting incumbent first completes the original 128-bin qualification, paid conventional repair, refinement ladder and exact fallback if needed. Only then do safeguarded variants try three original-support steps. Each accepted action must have a fresh policy upper bound strictly below the current policy lower bound and must pass both original targets. Refused trials are retained and charged.

\begin{table}[htbp]\caption{Complete-service clocks and standalone acquisition}\label{tab:clocks109}
\centering\small
\begin{tabular}{@{}lrrrrrr@{}}\toprule
&\multicolumn{3}{c}{$\alpha=1/2$}&\multicolumn{3}{c}{$\alpha=3/4$}\\
Representation&Acquire&Legacy&Compiled&Acquire&Legacy&Compiled\\\midrule
Online & 0.0001 & 0.5101 & 0.1020 & 0.0002 & 0.5262 & 0.1029\\
Quadratic (rank 6) & 0.0892 & 0.1139 & 0.0599 & 0.0856 & 0.1164 & 0.0591\\
ReLU (rank 6) & 0.0828 & 0.2490 & 0.0942 & 0.0735 & 0.2477 & 0.0958\\
Same-geometry hinge & 0.0830 & 0.2453 & 0.0826 & 0.0737 & 0.2470 & 0.0825\\
Coordinate spline & 0.0696 & 0.1120 & 0.0582 & 0.0703 & 0.1147 & 0.0578\\
Full quadratic & 0.0143 & 0.2766 & 0.2252 & 0.0146 & 0.2841 & 0.2217\\
Full cubic & 0.0155 & 0.3176 & 0.2549 & 0.0165 & 0.3118 & 0.2640\\
Bilinear grid & 0.0141 & 0.2888 & 0.2415 & 0.0144 & 0.2743 & 0.2270\\
\bottomrule\end{tabular}
\par\smallskip\begin{minipage}{.98\linewidth}\footnotesize
Seconds. Service entries are medians of all seven sequential fresh-process 32-state repetitions, including qualification and fallback. Acquisition includes paid common labels, fitting, independent numerical audit and serialization; it is charged in full to each standalone fitted arm. Hinge pays its shared ReLU acquisition. All clock ranges and cumulative process memory are in the Development Appendix. Full research, comparator selection, reference and timing work is separately recorded; these standalone charges are not the cost of selecting the best arm after seeing this table. Full polynomial and grid arms are strong resource-frontier controls, not equal-memory comparisons.
\end{minipage}\end{table}

\begin{table}[htbp]\caption{Actual policy costs and paid safeguarded improvement}\label{tab:economics109}
\centering\scriptsize
\begin{tabular}{@{}llrrr@{}}\toprule
$\alpha$&Representation&Compiled mean cost&Mean cost reduction&Changed states\\\midrule
$1/2$ & Online & [0.608374, 0.608376] & [0.00000000, 0.00000000] & 0/64\\
$1/2$ & Quadratic (rank 6) & [0.608522, 0.608525] & [0.00014797, 0.00014906] & 30/64\\
$1/2$ & ReLU (rank 6) & [0.608551, 0.608553] & [0.00017658, 0.00017720] & 17/64\\
$1/2$ & Same-geometry hinge & [0.608551, 0.608553] & [0.00017658, 0.00017720] & 17/64\\
$1/2$ & Coordinate spline & [0.608660, 0.608662] & [0.00028531, 0.00028606] & 21/64\\
$1/2$ & Full quadratic & [0.608522, 0.608525] & [0.00014798, 0.00014907] & 30/64\\
$1/2$ & Full cubic & [0.608588, 0.608590] & [0.00021357, 0.00021473] & 32/64\\
$1/2$ & Bilinear grid & [0.608389, 0.608391] & [0.00001461, 0.00001562] & 28/64\\
$3/4$ & Online & [0.618018, 0.618020] & [0.00000000, 0.00000000] & 0/64\\
$3/4$ & Quadratic (rank 6) & [0.618174, 0.618176] & [0.00015470, 0.00015578] & 30/64\\
$3/4$ & ReLU (rank 6) & [0.618172, 0.618174] & [0.00015302, 0.00015360] & 16/64\\
$3/4$ & Same-geometry hinge & [0.618172, 0.618174] & [0.00015302, 0.00015360] & 16/64\\
$3/4$ & Coordinate spline & [0.618307, 0.618309] & [0.00028786, 0.00028861] & 21/64\\
$3/4$ & Full quadratic & [0.618174, 0.618176] & [0.00015471, 0.00015579] & 30/64\\
$3/4$ & Full cubic & [0.618225, 0.618228] & [0.00020616, 0.00020730] & 32/64\\
$3/4$ & Bilinear grid & [0.618068, 0.618070] & [0.00004872, 0.00004970] & 27/64\\
\bottomrule\end{tabular}
\par\smallskip\begin{minipage}{.98\linewidth}\footnotesize
Outward numerical intervals average all 64 declared states, not a random sample. Cost reductions compare the improved policy with its completed same-backend incumbent and the same actual terminal policy. Identical actions have exactly zero contrast; every changed action passes strict interval separation. Every refused trial is charged. These gains do not assert that the extra computation is economical at an unmeasured compute price. Full safe-service clocks, both backends, sharpened regret and all pairwise original-cost contrasts are retained.
\end{minipage}\end{table}


All 4,096 catalogue decisions completed. Seven sequential fresh-process timing repetitions add 14,336 decisions, giving 18,432 retained decision records and 27,648 attempted safeguard records. The independent standard-library rational audit passes 169,472 stored-bound, feasibility, strict-gate and trace checks. Separately, 384 original-kernel contexts at $d=2,4,8$ pass 3,072 rational primitive, cubic-derivative and constrained-optimum checks. This dimension diagnostic is not a learned high-dimensional or longer-horizon tournament. Sixteen frozen-source preflight tests, fifteen retained exact-mechanism tests and twenty-three additional algebra/accounting regressions also pass; finite tests do not replace the continuum proofs.

The compiler's benefit is not limited to stored neural continuations. The online 32-state median falls from about $0.510$ to $0.102$ seconds at $\alpha=1/2$, and from $0.526$ to $0.103$ at $3/4$. Rank-six quadratic falls from $0.114$ to $0.060$ and from $0.116$ to $0.059$ seconds. The corresponding ReLU clocks are $0.249$ to $0.094$ and $0.248$ to $0.096$. These seven-repetition medians are descriptive host measurements; frequency and external load were not controlled. The full ranges, original per-category operation events and cumulative process memory are retained. Heterogeneous events are not summed and called FLOPs.

At $\alpha=1/2$, the unsafeguarded compiled ReLU's mean original cost lies in $[0.608551,0.608553]$, versus $[0.608522,0.608525]$ for rank-six quadratic and $[0.608374,0.608376]$ for online. At $3/4$, ReLU minus quadratic has a certified finite-catalogue mean difference in $[-0.00000241,-0.00000127]$. That small improvement reverses the low-risk ordering, does not beat the online policy, costs more in this implementation, and is exactly reproduced by conventional hinges. Thus neither selecting the favorable risk level nor renaming equivalent geometry establishes neural specificity.

Across both backends and risk levels, safeguards accept 1,340 update events; multiple updates can affect one state. For compiled ReLU, 17 of 64 states improve at risk $1/2$, with mean reduction in $[0.00017658,0.00017720]$; at $3/4$, 16 states improve, with mean reduction in $[0.00015302,0.00015360]$. These reductions compare actual policy costs to a completed same-backend incumbent. The associated median 32-state service clocks rise from approximately $0.094$ to $0.288$ seconds and from $0.096$ to $0.278$. Online safeguards accept no change. The improvement mechanism is therefore economically valid but not automatically worth buying: compute prices and deployment workloads would have to justify these additional costs. No such calibration is invented here.

The ReLU and hinge arms have exactly equal complete row and trial records in all eight risk-by-service comparisons. Their different clocks expose implementation cost, not a difference in function class, acquired information, policy quality or certification. The grid control has a smaller unsafeguarded policy loss than rank-six quadratic on both risk catalogues, but its proposal search and larger stored object are paid. No comparator is removed because it wins on cost, policy quality or a particular resource axis.

\subsection{Historical adverse evidence and verification limits}
The R97 acquisition campaign remains a separate historical result: all 36 validated residual spans decreased, whereas first-proposal rejections rose from 644 to 757. Quadratic rejections rose from 91 to 210; ReLU changed from 286 to 287 and spline from 267 to 260. Quadratic least squares was that campaign's observed complete-work winner. Its 1,800-cell sensitivity grid certified 156 adoptions over the online solver, 199 rejections and 1,445 unresolved cells, with \emph{zero} learned-service adoption over quadratic least squares. All twelve retained R96 simultaneous policy-cost intervals cross zero. These results are neither erased nor relabeled as successful training for the tail-aware objective introduced here.

The whole comparative campaign, including reference work and all seven timing repetitions, took about $224.66$ seconds internally on the recorded host; independently charged reference computation accounts for about $6.14$ seconds. The complete research comparison is not replaced by the cheaper cost of deploying one arm chosen ex ante. Labels, every rejected feature trial, fitting audits, serialization and every refused action improvement remain in the resource ledger. Source inspection and code-development labor are not measured CPU expenses and are not assigned fictitious zero economic cost.

A full inherited-test replay was also attempted, rather than inferred from old release claims. In the delivered R99 package layout, 435 test cases ran across 21 suites: ten suites passed, while the aggregate run recorded twenty errors and one failed assertion. The missing objects include old imported kernels and pinned historical evidence paths; the failed assertion explicitly requires a missing pinned cost file. These logs are preserved. This prevents a claim of an all-green full-history replay. The current R108-source campaign and its imported dependency closure did execute successfully, and the current mathematical and PDF verification are reported separately. Restoring every historical external dependency remains a publication-archive obligation, not a reason to invent old evidence.

The present evidence establishes original-policy improvements, an original-model computational improvement and a fairer conventional frontier. It does not establish an independently calibrated welfare purchase, a population ordering for newly fitted actors, a neural-exclusive competitive region, or a longer-horizon multivariate continuous-risk solution. The NBO research program continues to target those questions within the original decision problem; current proofs and computations are not substituted for them.


\section{Conclusion}\label{sec:conclusion109}
Neural Bellman Operators remain a construction-and-learning approach to original recursive economic decisions. Their continuation maps must be judged by feasible decisions, original-objective losses and complete paid service work, rather than by an internal residual score alone. The centered tail-switch theorem identifies a specific reason that directional fitting can fail under nonlinear risk. Original-law support and strict completed-incumbent gates supply a verifiable route from a proposed action to an actual cost improvement. The recursive gain account and precision bound make explicit which uniform hypotheses and expenditures are needed beyond a finite successful catalogue.

In the original investment economy, the active-region Bellman compiler provides a strong semi-analytic control and a useful common backend. The executed catalogue confirms strict policy improvements but also their additional cost. It confirms exact equivalence between acquired ReLU geometry and conventional hinges, not a uniquely neural advantage. Conventional victories, adverse residual-fitting results and unresolved adoption regions remain part of the scientific record.

The next scientific tests are substantive: prospectively train and evaluate the tail-aware objective under independently verified margin information; establish a neural representation advantage against equally informed conventional geometry and resources; solve genuinely longer-horizon, multiple-control and multiple-innovation instances; and calibrate where certified policy gains justify acquisition and service expense. The Brownian recursive-utility bridge states one rigorous connection to NDU/FBSDE research without transferring evidence across different objectives. Neither the finite comparison nor successful typesetting completes these broader claims. The original economic problem, mathematical ambition and numerical-methods objective are retained.

\bibliographystyle{ecta-fullname}
\bibliography{references99,references109-additions}
\end{document}
